\documentclass[12pt]{article}

% Idioma e codificação
\usepackage[brazil]{babel}
\usepackage[utf8]{inputenc}
\usepackage[T1]{fontenc}

% Matemática e listas
\usepackage{amsmath, amssymb}
\usepackage{enumitem}

% Layout e cores
\usepackage{geometry}
\usepackage{graphicx}
\usepackage{xcolor}
\usepackage{tikz}
\usetikzlibrary{positioning, arrows.meta}

% Cores institucionais
\definecolor{maincolor}{RGB}{0,102,51} % Verde UniRV
\definecolor{meuazul}{RGB}{0, 79, 151}

% Cabeçalho e rodapé
\usepackage{fancyhdr}

% Blocos
\usepackage[most]{tcolorbox}
\newtcolorbox{block}[1]{%
  colback=green!5,
  colframe=maincolor,
  coltitle=white,
  title=#1,
  fonttitle=\bfseries,
  boxrule=1.5pt,
  arc=3mm,
  left=6pt,
  right=6pt,
  top=6pt,
  bottom=6pt
}

\newtcolorbox{resposta}[1]{%
  colback=white,
  colframe=maincolor!30,
  arc=1mm,
  height=#1,
  width=\textwidth,
  boxrule=0.8pt,
  breakable
}

% Margens
\geometry{margin=2.5cm}

% Fancyhdr
\pagestyle{fancy}
\fancyhf{}

\renewcommand{\headrulewidth}{2pt}
\renewcommand{\footrulewidth}{2pt}

\renewcommand{\headrule}{\hbox to\headwidth{%
  \color{maincolor}\leaders\hrule height \headrulewidth\hfill}}

\renewcommand{\footrule}{\hbox to\headwidth{%
  \color{maincolor}\leaders\hrule height \footrulewidth\hfill}}

\setlength{\headheight}{52pt}
\addtolength{\topmargin}{-20pt}

\fancyhead[L]{%
  \textbf{UniRV} \\ \small Universidade de Rio Verde
}

\fancyhead[R]{%
  \raggedleft
  \textcolor{maincolor}{%
    \small
    \textbf{Lista de Revisão (Aulas 1 a 11)}\\
    Prof. Me.~Sergio Souza Novak\\
    \textit{Engenharia da Qualidade e Confiabilidade}
  }
}

\fancyfoot[C]{\textcolor{maincolor}{\thepage}}

% =========================
\begin{document}
% =========================

\begin{center}
  {\Large \textbf{Lista de Revisão Geral: Engenharia da Qualidade e Confiabilidade}}\\[0.2cm]
  \textit{Conteúdo abrangendo os fundamentos, métricas, monitoração e segurança.}
\end{center}

\vspace{0.4cm}

\begin{block}{Resumo de Fórmulas e Conceitos Chave}
  \small
  \begin{itemize}
    \item \textbf{POFOD}: $P(\text{falha na demanda}) = \frac{\text{Nº de Falhas}}{\text{Nº de Solicitações}}$
    \item \textbf{ROCOF}: $\text{Taxa de Ocorrência de Falhas} = \frac{\text{Nº de Falhas}}{\text{Tempo Total}}$
    \item \textbf{MTTF}: $1 / \text{ROCOF}$ (Tempo Médio para Falhas)
    \item \textbf{MTTR}: $\frac{\text{Tempo Total de Reparo}}{\text{Nº de Falhas}}$ (Tempo Médio para Reparo)
    \item \textbf{Disponibilidade (A)}: $\frac{\text{MTTF}}{\text{MTTF} + \text{MTTR}}$
    \item \textbf{4 Sinais de Ouro}: Latência, Tráfego, Erros e Saturação.
    \item \textbf{CIA Triad}: Confidencialidade, Integridade e Disponibilidade.
  \end{itemize}
\end{block}




\textbf{1.} Explique a principal diferença entre \textbf{Qualidade} (foco na entrega) e \textbf{Confiabilidade} (foco no tempo). Por que um sistema pode ser considerado de "alta qualidade" no dia do lançamento e apresentar "baixa confiabilidade" no mês seguinte?
\begin{resposta}{3.5cm}
\end{resposta}

\textbf{2.} Sobre o programa Apollo e o trabalho de \textbf{Margaret Hamilton}:
\begin{enumerate}[label=\alph*)]
    \item O que Hamilton defendia em relação ao "erro humano" no design de sistemas?
    \item Explique como o incidente do programa \textbf{P01} na Apollo 8 demonstrou a necessidade de sistemas resilientes (tolerância a falhas).
\end{enumerate}
\begin{resposta}{4.5cm}
\end{resposta}

\textbf{3.} Diferencie o modelo tradicional de \textbf{Sysadmin} (Administrador de Sistemas) do modelo de \textbf{SRE} (Site Reliability Engineering). Qual o impacto do crescimento do serviço no tamanho da equipe em cada modelo?
\begin{resposta}{4cm}
\end{resposta}




\textbf{4.} Um servidor de banco de dados operou por um período de 30 dias ($720$ horas). Durante este tempo, ocorreram $4$ interrupções. O tempo total para restaurar o serviço em todas as interrupções somadas foi de $3$ horas.
\begin{enumerate}[label=\alph*)]
    \item Calcule o \textbf{MTTF} e o \textbf{MTTR}.
    \item Determine a \textbf{Disponibilidade} percentual do sistema.
    \item Se o requisito do cliente é "três noves" ($99,9\%$), o sistema está em conformidade?
\end{enumerate}
\begin{resposta}{6cm}
\end{resposta}

\textbf{5.} Compare os dois sistemas abaixo:
\begin{itemize}
    \item \textbf{Sistema X}: Falha a cada 1.000h, mas leva 10h para ser reparado manualmente.
    \item \textbf{Sistema Y}: Falha a cada 100h, mas possui recuperação automática que leva apenas 30 segundos (0,0083h).
\end{itemize}
Calcule a disponibilidade de ambos e discuta por que, em SRE, muitas vezes prefere-se o sistema que falha mais, porém recupera-se muito mais rápido.
\begin{resposta}{5.5cm}
\end{resposta}




\textbf{6.} O Google define os \textbf{"4 Sinais de Ouro"} da monitoração. Liste-os e explique brevemente como cada um ajuda a diagnosticar a saúde de um sistema web.
\begin{resposta}{4.5cm}
\end{resposta}

\textbf{7.} Defina os conceitos de \textbf{SLI}, \textbf{SLO} e \textbf{Error Budget}. Como o Error Budget (Orçamento de Erro) ajuda a resolver o conflito entre o time de Desenvolvimento (que quer velocidade) e o time de Operações (que quer estabilidade)?
\begin{resposta}{4.5cm}
\end{resposta}

\textbf{8.} Por que a \textbf{média aritmética} de latência pode ser uma métrica perigosa em sistemas distribuídos? Qual a vantagem de utilizar \textbf{Percentis} (ex: P95, P99) para medir a experiência do usuário?
\begin{resposta}{4cm}
\end{resposta}




\textbf{9.} O que é um \textbf{Postmortem "Sem Culpa"} (\textit{Blameless Postmortem})? Por que o Google e outras empresas de tecnologia evitam punir indivíduos por falhas sistêmicas?
\begin{resposta}{4cm}
\end{resposta}

\textbf{10.} Explique os conceitos de:
\begin{enumerate}[label=\alph*)]
    \item \textbf{Runbook}: Qual sua utilidade durante um incidente às 3h da manhã?
    \item \textbf{Canary Deployment}: Como essa estratégia reduz o "raio de explosão" de um bug em produção?
\end{enumerate}
\begin{resposta}{5cm}
\end{resposta}




\textbf{11.} Explique a tríade \textbf{CIA} (Confidencialidade, Integridade e Disponibilidade). Dê um exemplo de um ataque que viole cada um desses pilares.
\begin{resposta}{4.5cm}
\end{resposta}

\textbf{12.} Sobre Criptografia:
\begin{enumerate}[label=\alph*)]
    \item Diferencie Criptografia \textbf{Simétrica} de \textbf{Assimétrica}.
    \item Por que usamos Criptografia de Chave Pública (Assimétrica) para trocar chaves de sessão, mas usamos Criptografia Simétrica (como AES) para o tráfego de dados volumosos?
\end{enumerate}
\begin{resposta}{5cm}
\end{resposta}

\textbf{13.} O que é o \textbf{Não-Repúdio} na segurança da informação e como as Assinaturas Digitais garantem esse atributo?
\begin{resposta}{3.5cm}
\end{resposta}

\vspace{1cm}
\begin{center}
    \textit{“A sorte favorece a mente preparada. Boa revisão!”} \\
    \textbf{Prof. Me. Sergio Souza Novak}
\end{center}

\end{document}
